\definecolor{asioblue}{RGB}{63,94,143}
\definecolor{bridgegreen}{RGB}{69,126,102}
\definecolor{kernelyellow}{RGB}{173,136,55}
\colorlet{fiberfill}{idlekvred!7}
\colorlet{asiofill}{asioblue!7}
\colorlet{bridgefill}{bridgegreen!8}
\colorlet{kernelfill}{kernelyellow!10}

\begin{tikzpicture}[
  node distance=0.38cm,
  stage/.style={
    draw=#1,
    fill=#1!8,
    rounded corners=2pt,
    line width=0.8pt,
    minimum width=11.4cm,
    minimum height=0.95cm,
    align=left,
    inner xsep=8pt,
    inner ysep=6pt,
    font=\small
  },
  syslabel/.style={
    font=\small\bfseries,
    text width=2.2cm,
    align=right
  },
  flow/.style={
    -{Latex[length=2.6mm,width=1.8mm]},
    line width=0.9pt,
    draw=black!75
  }
]

\node[syslabel] (l1) at (0,0) {Fiber};
\node[stage=idlekvred, anchor=west] (n1) at (2.2,0)
  {连接 Fiber 执行命令逻辑，并调用 \texttt{async\_read(..., yield\_t)}}; 

\node[syslabel, below=of l1] (l2) {桥接层};
\node[stage=bridgegreen, anchor=west, below=of n1] (n2)
  {completion token 桥接层扩展 \texttt{yield\_t} 与 \texttt{async\_result}，保存当前 Fiber 的恢复点}; 

\node[syslabel, below=of l2] (l3) {Asio};
\node[stage=asioblue, anchor=west, below=of n2] (n3)
  {\texttt{io\_context} 为 socket 注册异步读写操作，并将完成处理器关联到事件循环}; 

\node[syslabel, below=of l3] (l4) {EventLoop};
\node[stage=asioblue, anchor=west, below=of n3] (n4)
  {当前 Fiber 挂起，工作线程返回事件循环，继续调度其他 ready Fiber}; 

\node[syslabel, below=of l4] (l5) {Kernel};
\node[stage=kernelyellow, anchor=west, below=of n4] (n5)
  {socket 变为可读/可写，内核将 I/O 完成事件返回给事件循环}; 

\node[syslabel, below=of l5] (l6) {Asio};
\node[stage=asioblue, anchor=west, below=of n5] (n6)
  {Asio 在同一工作线程中执行完成回调，并把结果交回 completion token}; 

\node[syslabel, below=of l6] (l7) {桥接层};
\node[stage=bridgegreen, anchor=west, below=of n6] (n7)
  {桥接层唤醒原 Fiber，回填返回值或错误码，恢复挂起的执行上下文}; 

\node[syslabel, below=of l7] (l8) {Fiber};
\node[stage=idlekvred, anchor=west, below=of n7] (n8)
  {原 Fiber 从挂起点继续执行，随后进入请求解析、命令执行或响应发送路径}; 

\draw[flow] (n1.south) -- (n2.north);
\draw[flow] (n2.south) -- (n3.north);
\draw[flow] (n3.south) -- (n4.north);
\draw[flow] (n4.south) -- (n5.north);
\draw[flow] (n5.south) -- (n6.north);
\draw[flow] (n6.south) -- (n7.north);
\draw[flow] (n7.south) -- (n8.north);

\end{tikzpicture}
